% TOP_PROBLEM: {"id": 1165, "title": "The Total Coloring Conjecture", "category_id": 3, "category": {"id": 3, "name": "graph_theory", "display_name": "Graph Theory", "description": "Problems involving graphs, networks, and their properties.", "slug": "graph-theory", "order_index": 3, "created_at": "2026-07-31T15:26:25.671Z"}, "status": "open"}
% FIRST_POSTED: 2026-10-01
% ATTEMPT_STATUS: unresolved
% Imported from: unsolvedmath_additional_500.tex; UnsolvedMath ID 1165
% !TeX program = lualatex
% Compile twice: lualatex 1165.tex
% Self-contained: no external bibliography, images, or \input files.
% Numeric section labels and visible numbers are original UnsolvedMath IDs.
\documentclass[11pt,a4paper]{article}
\usepackage[margin=24mm,headheight=14pt]{geometry}
\usepackage{fontspec}
\setmainfont{Latin Modern Roman}
\setsansfont{Latin Modern Sans}
\setmonofont{Latin Modern Mono}
\usepackage{amsmath,amssymb,mathtools}
\usepackage{unicode-math}
\setmathfont{Latin Modern Math}
\usepackage{microtype}
\usepackage{enumitem,xurl,needspace}
\usepackage[dvipsnames]{xcolor}
\usepackage{hyperref}
\hypersetup{unicode=true,colorlinks=true,linkcolor=MidnightBlue,urlcolor=MidnightBlue,
 pdftitle={UnsolvedMath Problem 1165},pdfauthor={Source-annotated compilation},bookmarksnumbered=true}
\usepackage{bookmark,tocloft,titlesec,fancyhdr}
\setlength{\cftsecnumwidth}{4.2em}
\cftsetpnumwidth{2.5em}
\cftsetrmarg{4em}
\titleformat{\section}{\Large\bfseries\raggedright}{\thesection}{0.7em}{}
\pagestyle{fancy}\fancyhf{}
\fancyhead[L]{\small\sffamily Additional UnsolvedMath statements}
\fancyhead[R]{\small Original catalogue IDs}
\fancyfoot[C]{\thepage}
\setlength{\parindent}{0pt}
\setlength{\parskip}{5pt plus 1pt}
\setlength{\emergencystretch}{3em}
\setcounter{tocdepth}{1}\setcounter{secnumdepth}{1}
\setlist[itemize]{leftmargin=3em,itemsep=4pt,topsep=4pt,parsep=0pt}
\allowdisplaybreaks
\newcommand{\N}{\mathbb{N}}
\newcommand{\Z}{\mathbb{Z}}
\newcommand{\Q}{\mathbb{Q}}
\newcommand{\R}{\mathbb{R}}
\newcommand{\C}{\mathbb{C}}
\newcommand{\F}{\mathbb{F}}
\newcommand{\A}{\mathbb{A}}
\newcommand{\PP}{\mathbb{P}}
\newcommand{\E}{\mathbb{E}}
\newcommand{\eps}{\varepsilon}
\newcommand{\dd}{\,\mathrm d}
\newcommand{\sref}[2]{\hyperlink{source:#1:#2}{[S#2]}}
\newif\ifshowcatalogue
\showcataloguetrue
\newcommand{\cataloguescope}[1]{\ifshowcatalogue
 \paragraph{Statement recorded in the supplied source.}
 #1\fi}
\begin{document}
% UnsolvedMath ID 1165; source catalogue code GT-010

\setcounter{section}{1164}

\section{The Total Colouring Conjecture}\label{1165}

{\small\sffamily UnsolvedMath ID 1165 · Graph Theory}

\subsection{Definitions and mathematical statement}

\paragraph{Shared notation.}Unless a section says otherwise, $\N=\{1,2,\ldots\}$,
$\Z$ is the ring of integers, $\Q,\R,\C$ are the rational, real, and complex
fields, $[N]=\{1,\ldots,\lfloor N\rfloor\}$, and $\log$ is the natural
logarithm. For real functions of a parameter tending to infinity,
$f\ll g$ and $f=O(g)$ mean $|f|\leq Cg$ eventually for some fixed $C$;
$f\gg g$ reverses this comparison; $f=o(g)$ means $f/g\to0$;
$f\sim g$ means $f/g\to1$; and $f\asymp g$ means both order bounds.
A subscript on $O$ or $\ll$ specifies allowed constant dependence.
The expression $n^{o(1)}$ in an upper bound means growth smaller than
$n^\epsilon$ for each fixed $\epsilon>0$. Local definitions govern any
exception, finite-field convention, or use of the same symbol for another
object. An asymptotic claim is not automatically an effective algorithm.



A total colouring assigns colours to both $V(G)$ and $E(G)$. Adjacent vertices, edges sharing a vertex, and an edge together with either endpoint must receive different colours. Write $\chi''(G)$ for the minimum number of colours and $\Delta(G)=\max_v\deg(v)$. The conjecture is $\chi''(G)\leq\Delta(G)+2$ for every finite simple graph. \sref{1165}{1} \sref{1165}{2}

\paragraph{Statement recorded in the supplied source.}
Can every graph be totally colored with at most $\Delta + 2$ colors, where $\Delta$ is the maximum degree? \sref{1165}{1}

\paragraph{Residual task reported by the archive.}

The embedded review (2026-08-17) records: Prove the Delta+2 total-coloring bound for every graph. \sref{1165}{1}

\subsection{Short English statement}

Can the vertices and edges of every graph be coloured without any incidence conflicts using at most two more colours than its maximum degree?

\subsection{Research attempt: constructive tree colourings and sharp complete-graph tests}
\textbf{Status and source scope.} This is a partial attempt, not a proof of the general conjecture. The inspected abstract of [S2] concerns a restricted planar class. The elementary constructions below are proved here and make no novelty claim.

\subsubsection{1. Reformulating the conflict constraints}
Form the total graph $T(G)$, whose vertices are the objects $V(G)\sqcup E(G)$ and whose adjacency is exactly a forbidden equality of colours in the question. A vertex-object $v$ has degree $2d_G(v)$; an edge-object $uv$ has degree $d_G(u)+d_G(v)$. Consequently greedy colouring proves
\[\Delta+1\leq\chi''(G)\leq2\Delta+1\]
when $G$ has an edge: the lower bound is the clique consisting of a maximum-degree vertex and its incident edges. For $\Delta=1$, a component $K_2$ actually forms a triangle in the total graph and needs three colours. This exposes a boundary exception to any proposed universal $\Delta+1$ argument.

\subsubsection{2. A complete construction for trees}
Let $T$ be a tree with maximum degree $\Delta\geq2$, and use a fixed palette of $\Delta+1$ colours. Root the tree. Colour the root arbitrarily and colour its incident edges distinctly with the other colours. Suppose a non-root vertex $v$ is reached after its parent $u$ and the parent edge $uv$ have been coloured. Choose the colour of $v$ to avoid the two distinct colours on $u$ and $uv$; there are at least $\Delta-1\geq1$ choices. There are then $\Delta-1$ colours other than the colours of $v$ and $uv$. Assign distinct members of this set to the at most $\Delta-1$ child edges of $v$.

Continue by distance from the root. Every vertex differs from its parent; every edge differs from both endpoints once the child is coloured; and the edges incident at each vertex have distinct colours. There are no additional conflicts, because a tree has no edges between different branches. Thus
\[\chi''(T)=\Delta+1\qquad(\Delta\geq2).\]
The argument also handles forests componentwise, with isolated edges requiring three colours and isolated vertices requiring one. It is constructive and does not need an edge-colouring theorem.

\subsubsection{3. Why the extra colour can be necessary}
Let $n\geq2$ be even and consider $K_n$. A total colouring with $n$ colours would give its $n$ vertices all different colours, so every colour occurs on exactly one vertex. Edges of a fixed colour form a matching avoiding that vertex, and hence contain at most $(n-2)/2$ edges. Across all $n$ colours this accounts for at most $n(n-2)/2$ edges, fewer than the actual $n(n-1)/2$. Therefore $\chi''(K_n)\geq n+1=\Delta+2$.

For the reverse inequality put $q=n+1$, an odd integer, label $n$ vertices by distinct residues modulo $q$, colour vertex $i$ by $2i$, and colour edge $ij$ by $i+j$. Doubling is injective modulo odd $q$. At a fixed vertex $i$, the edge colours are distinct and none equals $2i$; the colour of $ij$ differs from $2j$ as well. This proves the exact value $n+1$. The same construction on all $q$ vertices gives $\chi''(K_q)=q$ for odd $q\geq3$.

\subsubsection{4. Adversarial check and the unresolved step}
The tree construction fails on a graph with cycles because a newly coloured vertex can already have several coloured neighbours, and a newly processed edge can impose constraints at both ends. Treating those extra conflicts as absent would be an invalid proof. Conversely, the complete-graph test rules out replacing the target by $\Delta+1$. The accompanying diagnostic checks every incidence in the modular construction for complete graphs of orders 2 through 30 and checks the tree construction on every labelled tree through six vertices. These are implementation checks of the constructions, not evidence covering all graphs. A method that coordinates the two endpoint constraints using only $\Delta+2$ colours in an arbitrary graph is still missing.

\subsection{Sources}

{\small

\begin{itemize}

\item[\hypertarget{source:1165:1}{S1}] ULAM.AI, \emph{UnsolvedMath}, supplied \texttt{problems.json}, record 1165 (GT-010), statement and background. The embedded literature-review date is 2026-08-17; that is the archive’s review date, not a fresh certification. Dataset landing page: \url{https://huggingface.co/datasets/ulamai/UnsolvedMath}.

\item[\hypertarget{source:1165:2}{S2}] R. Su, G. Fang and E. Zhu, The Total Coloring Conjecture holds for planar graphs without three special subgraphs, arXiv:2507.12737 (2025). \url{https://arxiv.org/abs/2507.12737}. Bibliographic information imported from the supplied record.

\end{itemize}

}

\label{end:1165}
\end{document}
